% Copyright (c) 2026 Selina Cheggour. Tous droits réservés. Voir le fichier LICENSE.
% Annexe : les méthodes eom, leaf (HDBSCAN) et Ward (plan B), exercices 3.1 et 3.3.

\newpage
\section{Trois façons de faire les tas : eom, leaf et Ward (exercices 3.1 et 3.3)}

\subsection{D'abord, l'arbre de HDBSCAN}

HDBSCAN regarde les plaintes comme des îles sur une carte, et fait \textbf{baisser le niveau de la mer}
petit à petit. Au début, tout est une seule grande île. Quand l'eau monte entre deux zones moins denses,
l'île se coupe en deux, puis encore en deux\ldots{} On obtient un \textbf{arbre} : chaque branche est un
groupe possible, et sa \textbf{longueur} dit combien de temps il a tenu avant de se couper. Une branche de
moins de \texttt{taille\_min} plaintes ne compte pas : ses plaintes deviennent du bruit (étiquette $-1$).

Il reste à choisir \textbf{quelles branches garder}. C'est là qu'interviennent les deux méthodes.

\begin{center}
\begin{tikzpicture}[x=0.62cm, y=0.62cm,
  br/.style={line width=#1, color=slate!70},
  pris/.style={line width=#1, color=accentblue},
  lab/.style={font=\scriptsize, color=slate}]
\foreach \dx/\titre/\coul in {0/{\texttt{eom} : les branches qui durent}/accentblue, 12.5/{\texttt{leaf} : les bouts de l'arbre}/ambre} {
\begin{scope}[xshift=\dx cm]
  \node[font=\small\bfseries, color=\coul] at (5,8.4) {\titre};
  % tronc
  \draw[br=5pt] (5,7.6) -- (5,6.5);
  \draw[br=2pt] (2.5,6.5) -- (7.5,6.5);
  % branche gauche A : longue, puis coupée en deux petits bouts
  \draw[br=4pt] (2.5,6.5) -- (2.5,2.6);
  \draw[br=1.5pt] (1.5,2.6) -- (3.5,2.6);
  \draw[br=2pt] (1.5,2.6) -- (1.5,1.4);
  \draw[br=2pt] (3.5,2.6) -- (3.5,1.6);
  % branche droite B : courte, puis coupée en deux longues branches
  \draw[br=4pt] (7.5,6.5) -- (7.5,5.2);
  \draw[br=1.5pt] (6.3,5.2) -- (8.7,5.2);
  \draw[br=3pt] (6.3,5.2) -- (6.3,0.8);
  \draw[br=3pt] (8.7,5.2) -- (8.7,1.2);
  \node[lab] at (2.5,7.0) {A}; \node[lab] at (7.5,7.0) {B};
  \node[lab, below] at (1.5,1.4) {A1}; \node[lab, below] at (3.5,1.6) {A2};
  \node[lab, below] at (6.3,0.8) {B1}; \node[lab, below] at (8.7,1.2) {B2};
\end{scope}
}
% eom : A (longue) et B1, B2 (longues) ; leaf : les quatre bouts
\draw[pris=4pt, color=accentblue] (2.5,6.5) -- (2.5,2.6);
\draw[pris=3pt, color=accentblue] (6.3,5.2) -- (6.3,0.8);
\draw[pris=3pt, color=accentblue] (8.7,5.2) -- (8.7,1.2);
\begin{scope}[xshift=12.5cm]
\draw[pris=2pt, color=ambre] (1.5,2.6) -- (1.5,1.4);
\draw[pris=2pt, color=ambre] (3.5,2.6) -- (3.5,1.6);
\draw[pris=3pt, color=ambre] (6.3,5.2) -- (6.3,0.8);
\draw[pris=3pt, color=ambre] (8.7,5.2) -- (8.7,1.2);
\end{scope}
\node[font=\footnotesize, text width=6.6cm, align=center] at (5,-0.6) {3 groupes : A entier, B1, B2};
\node[font=\footnotesize, text width=6.6cm, align=center] at (17.5,-0.6) {4 groupes : A1, A2, B1, B2};
\end{tikzpicture}
\end{center}
{\small\color{slate}\textit{Arbre inventé. En couleur, les branches gardées : chaque plainte ne peut être que dans une seule.}}

\begin{tcbraster}[raster columns=2, raster equal height, raster column skip=8pt, raster every box/.style={unbreakable}]
\begin{conceptbox}[title={\faTree\ \texttt{eom} : \og excess of mass \fg}, unbreakable]
\small Garde les branches qui \textbf{durent le plus longtemps}, en tenant compte de leur taille. Si une
branche dure plus que ses deux sous-branches réunies, on la garde entière (A) ; sinon on garde les
sous-branches (B1 et B2). Résultat : \textbf{moins de groupes, plus gros et plus solides}. C'est le choix
par défaut de HDBSCAN.
\end{conceptbox}
\begin{conceptbox}[title={\faLeaf\ \texttt{leaf} : les feuilles}, unbreakable]
\small Garde tous les \textbf{bouts de l'arbre}, les plus petites branches qui ne se coupent plus.
Résultat : \textbf{plus de groupes, plus petits et plus homogènes}. Utile quand on cherche des
problèmes fins, au risque de couper un même problème en morceaux.
\end{conceptbox}
\end{tcbraster}

\image{une famille qui partage un gâteau : \texttt{eom} coupe les grosses parts qui ont du sens (le côté
chocolat, le côté fraise), \texttt{leaf} coupe en toutes petites parts.}

\subsection{Le plan B : Ward}

Si aucun des 18 réglages de HDBSCAN n'est admis, c'est que les plaintes ne forment pas d'îles nettes.
On change alors de méthode : la classification de \textbf{Ward} part dans l'autre sens. Au début, chaque
plainte est seule dans son groupe ; à chaque étape, on \textbf{colle les deux groupes} dont la fusion
mélange le moins les plaintes, et on s'arrête quand il reste $k$ groupes.

\noindent\begin{minipage}{\linewidth}
\begin{center}
\begin{tikzpicture}[x=0.95cm, y=0.95cm, pt/.style={circle, fill=#1, inner sep=2.2pt},
  fl/.style={-{Stealth[length=2.5mm]}, thick, color=slate}]
\foreach \dx/\titre in {0/{départ : 6 groupes}, 4.6/{étape suivante : 4}, 9.2/{arrêt à $k = 2$}} {
  \node[font=\footnotesize\bfseries, color=darkblue] at (\dx+1.5,2.6) {\titre};
  \foreach \x/\y in {0.2/0.4, 0.8/0.9, 0.5/1.6, 2.2/0.5, 2.8/1.1, 2.4/1.8} \node[pt=sectioncolor] at (\dx+\x,\y) {};
}
% étape 2 : deux paires collées
\draw[accentblue, thick, rounded corners=6pt] (4.6+0.0,0.15) rectangle (4.6+1.05,1.15);
\draw[accentblue, thick, rounded corners=6pt] (4.6+1.95,0.25) rectangle (4.6+3.05,1.35);
% étape 3 : deux groupes
\draw[accentblue, very thick, rounded corners=8pt] (9.2-0.05,0.1) rectangle (9.2+0.85,1.9);
\draw[accentblue, very thick, rounded corners=8pt] (9.2+1.95,0.2) rectangle (9.2+3.05,2.1);
\draw[fl] (3.2,1.0) -- (4.3,1.0);
\draw[fl] (7.8,1.0) -- (8.9,1.0);
\end{tikzpicture}
\end{center}
{\small\color{slate}\textit{Exemple inventé avec 6 plaintes : à chaque étape, on colle les deux groupes les plus proches au sens de Ward.}}
\end{minipage}


\begin{formulabox}
\og Mélanger le moins \fg{} veut dire faire grandir le moins possible la \textbf{dispersion à l'intérieur des
groupes} :
\[
W = \sum_{h} \sum_{i \in h} \lVert \mathbf{z}_i - \bar{\mathbf{z}}_h \rVert^2 ,
\]
où $h$ parcourt les groupes, $\mathbf{z}_i$ est le vecteur de la plainte $i$ et $\bar{\mathbf{z}}_h$ la
moyenne du groupe $h$. Ward choisit à chaque étape la fusion qui augmente $W$ le moins.
\end{formulabox}

\begin{center}
\renewcommand{\arraystretch}{1.3}
\rowcolors{2}{notebg}{white}
\begin{tabular}{@{}>{\raggedright\arraybackslash}p{4.2cm}>{\raggedright\arraybackslash}p{5cm}>{\raggedright\arraybackslash}p{5cm}@{}}
\rowcolor{titlebg} & \color{white}\textbf{HDBSCAN (\texttt{eom} ou \texttt{leaf})} & \color{white}\textbf{Ward (plan B)} \\
nombre de groupes & trouvé tout seul & fixé d'avance : $k = 8, 12, 16, 20$ ou $30$ \\
plaintes isolées & laissées de côté ($-1$) & toujours rangées dans un groupe \\
ce qu'il cherche & des zones denses, de n'importe quelle forme & des groupes compacts, plutôt ronds \\
quand l'utiliser & des îles nettes dans les données & pas d'îles nettes : un continuum \\
\end{tabular}
\end{center}

\begin{piege}
\textbf{Ward range tout, même ce qui ne ressemble à rien.} Il n'y a plus de plaintes de côté : une
plainte bizarre finit quand même dans un groupe. C'est pourquoi on garde les mêmes règles d'admission
(pas de géant, assez de problèmes) et la fusion des jumeaux après Ward.
\end{piege}

\subsection{HDBSCAN et Ward : deux arbres, une seule différence qui compte}

Les deux méthodes sont \textbf{hiérarchiques}. Toutes les deux construisent un arbre de groupes emboîtés,
HDBSCAN en regardant la mer monter, Ward en collant les groupes deux à deux, puis décident où couper cet
arbre. Mais \textbf{seul HDBSCAN regarde la densité}. Ward ne la regarde jamais : il regarde la
\textbf{dispersion}. À chaque étape, il colle les deux groupes dont la fusion fait le moins grossir la
dispersion totale $W$, et il s'arrête quand il reste $k$ groupes.

\begin{tcbraster}[raster columns=2, raster equal height, raster column skip=8pt, raster every box/.style={unbreakable}]
\begin{conceptbox}[title={\faWater\ HDBSCAN}, unbreakable]
\small Un point isolé au milieu de nulle part, HDBSCAN le \textbf{laisse de côté} ($-1$), parce qu'il
n'est pas dans une zone dense. Il cherche \textbf{où c'est serré} et accepte de ne pas tout ranger.
\end{conceptbox}
\begin{conceptbox}[title={\faObjectGroup\ Ward}, unbreakable]
\small Le même point isolé, Ward le \textbf{colle de force} au groupe le moins loin, parce qu'il doit
tout ranger. Il \textbf{range tout} et cherche des paquets ronds et équilibrés.
\end{conceptbox}
\end{tcbraster}

\begin{lire}[unbreakable]
\textbf{Et pour les incidents ?} Un incident est une \textbf{zone dense} au milieu de plaintes variées :
c'est pour ça que HDBSCAN est le premier choix. Ward n'est que le plan B, quand les plaintes ne forment pas
d'îles nettes mais un \textbf{continuum}. Là, une méthode de densité ne trouve rien, ou un seul bloc, alors
que Ward découpe quand même quelque chose d'exploitable, avec les mêmes règles d'admission derrière.\par\smallskip
Donc Ward n'est pas \og ajouté après HDBSCAN \fg{} : il est essayé après lui, \textbf{à sa place},
uniquement si HDBSCAN a échoué sur toute la grille. Dans l'exécution normale, \textbf{Ward ne tourne pas
du tout}.
\end{lire}
